\newpage
\setcounter{table}{0}
\setcounter{figure}{0}
\section{系统实现及测试}

\subsection{硬件与软件环境}
\label{subsec:hardware_software}

实验集群共包含三台物理服务器，分别承担源端数据库（Source DB）、回放控制端（Replayer）和目标端数据库（Target DB）的角色。为了消除虚拟化技术带来的资源调度开销及不可预测的性能抖动，所有组件均部署在裸金属服务器（Bare Metal Server）上。各节点的详细硬件配置如表~\ref{tab:hardware_config} 所示。

% TODO: 这里应该替换为实际的机器配置
\begin{table}[htbp]
    \centering
    \caption{实验环境硬件配置表}
    \label{tab:hardware_config}
    \renewcommand{\arraystretch}{1.2}
    \begin{tabular}{lccc}
        \toprule
        \textbf{组件角色} & \textbf{CPU (Intel Xeon)} & \textbf{内存 (DDR4)} & \textbf{存储 (NVMe SSD)} \\
        \midrule
        Source DB  & Gold 6248R (24C/48T) & 128 GB & 2 TB (RAID 0) \\
        Replayer  & Silver 4210R (10C/20T) & 64 GB & 1 TB \\
        Target DB  & Gold 6248R (24C/48T) & 128 GB & 2 TB (RAID 0) \\
        \bottomrule
    \end{tabular}
\end{table}

在软件环境方面，操作系统统一采用 CentOS Linux 7.9 (Kernel 5.4.x)。数据库管理系统（Database Management System, DBMS）选用 PostgreSQL 18.1 版本，并开启 standard\_conforming\_strings 等标准配置以配合审计插件。本文设计的回放系统基于 Go 1.21 编译构建，利用其 Goroutine 调度机制实现高并发流量的解析与重放。源端安装并配置了 \texttt{pgaudit} 扩展插件，用于生成审计日志流。

\subsubsection{网络环境构建}

鉴于第 4.3.2 节中提出的基于时间戳的事务时序对齐算法对网络延迟具有极高的敏感度，实验环境的网络构建是本研究的关键环节。为了避免公网或共享局域网中的背景流量（Background Traffic）干扰，三台服务器通过万兆以太网交换机（10GbE Switch）构建了独立的物理局域网。
% TODO 机器直接的网络配置也应该替换为实际数据
各节点均配备双路 10Gbps 网络接口控制器（Network Interface Controller, NIC），并开启了网卡多队列（RSS）与直接内存访问（DMA）功能，以降低 CPU 的软中断开销。经 iPerf3 工具实测，Replayer 与 Target DB 之间的链路带宽稳定在 9.4 Gbps 以上，平均往返时延（Round-Trip Time, RTT）控制在微秒级，具体测试数据如下公式~\ref{eq:network_latency} 所示：

\begin{equation}
\label{eq:network_latency}
RTT_{avg} \approx 0.12ms, \quad \sigma_{RTT} < 0.05ms
\end{equation}

其中 $\sigma_{RTT}$ 表示时延的标准差。这种低延迟、低抖动的确定性网络环境，能够有效隔离网络传输因素，确保回放测试中观察到的时序偏差主要源于数据库内部的锁竞争或系统调度，而非网络传输瓶颈，从而验证回放引擎本身的时序控制能力。

\subsection{系统部署拓扑}
\label{subsec:deployment_topology}

流量回放系统的整体部署拓扑如图~\ref{fig:topology} 所示。整个数据链路严格遵循“生产-采集-回放-验证”的单向流动逻辑，确保测试过程不会对源端生产环境产生任何侵入性影响。

\begin{figure}[htbp]
    \centering
    \includegraphics[width=0.9\textwidth]{images/topology.pdf} 
    \caption{流量回放系统实验部署拓扑图}
    \label{fig:topology}
\end{figure}

具体的数据流转流程描述如下：

\begin{enumerate}
    \item \textbf{流量生成与捕获}：Source DB 承载由 Sysbench \cite{sysbench} 工具生成的模拟 OLTP 负载。数据库内核集成的审计模块实时捕获 SQL 执行记录，并将其序列化为 CSV 格式的审计日志文件。
    
    \item \textbf{日志传输与解析}：为了解耦生产库与回放端的计算压力，审计日志通过异步文件传输协议实时同步至 Replayer 节点。Replayer 内部的解析器（Parser）模块读取日志流，重构出包含事务上下文（Transaction Context）的中间表示结构。
    
    \item \textbf{并发回放}：Replayer 节点依据日志中的时间戳信息，通过维护一个全局虚拟时钟（Virtual Clock）来调度并发请求。它作为 Target DB 的客户端，建立连接池并复现原始流量的并发特征。
    
    \item \textbf{结果收集}：Target DB 在执行回放请求的同时，记录自身的性能指标（如 TPS, QPS, Latency），用于后续与 Source DB 的基准数据进行一致性比对。
\end{enumerate}

在该拓扑中，Replayer 节点部署于 Source DB 与 Target DB 之间，充当了流量整形网关（Traffic Shaping Gateway）的角色。物理上，Replayer 与 Target DB 的物理距离被刻意缩短（同机架部署），以模拟“本地回放”的效果，最大程度减少网络传输时延对高并发压力测试结果的衰减。
\subsection{核心模块实现演示}
\subsection{功能测试}
\subsection{性能与回放保真度评估}
\subsection{本章小结}

XXXXX